\section*{अध्याय \ref{ch:b2:anthropogenic-climate-change} --- मानवजनित जलवायु परिवर्तन}

\begin{solution}{exo:b2:anthropogenic-climate-change:1}
(क) $3.75 - 0.60 + 0.22 + 0.45 + 0.11 = \qty{3.93}{W.m^{-2}}$, समायोजन मिलकर
\qty{0.18}{W.m^{-2}} जोड़ते हैं। (ख) $3.93/3.71 = 1.059$: \qty{5.9}{\percent} अधिक। (ग) क्षोभमंडलीय
तापमान समायोजन, $-0.60$: गर्म क्षोभमंडल अंतरिक्ष में अधिक अवरक्त उत्सर्जित करता है, जो सतह के
अनुक्रिया देने से पहले ही अतिरिक्त ऊर्जा का कुछ भाग हटा देता है।
\end{solution}

\begin{solution}{exo:b2:anthropogenic-climate-change:2}
(क) $4\pi(\num{6.371e6})^2 = \qty{5.10e14}{m^2}$। (ख) 47 वर्षों पर, $\num{434.9e21}/(47 \times
\num{3.156e7} \times \num{5.10e14}) = \qty{0.575}{W.m^{-2}}$; 12 वर्षों पर, $\num{152.4e21}/(12
\times \num{3.156e7} \times \num{5.10e14}) = \qty{0.789}{W.m^{-2}}$, मूल्यांकन के 0.57 और
0.79। (ग) $0.91 \times 434.9 = \qty{396}{ZJ}$।
\end{solution}

\begin{solution}{exo:b2:anthropogenic-climate-change:3}
(क) $-3.22 + 1.30 + 0.35 + 0.42 - 0.01 = \num{-1.16}$, शुद्ध मान। (ख) $3.93/1.16 =
\qty{3.39}{\celsius}$; $3.93/3.22 = \qty{1.22}{\celsius}$। (ग) $3.39/1.22 = 2.78$: प्रतिपुष्टियाँ
प्लांक तापन को लगभग तिगुना कर देती हैं।
\end{solution}

\begin{solution}{exo:b2:anthropogenic-climate-change:4}
(क) एक पेटाग्राम कार्बन $44.009/12.011 = 3.664$~\unit{Gt} कार्बन डाइऑक्साइड है: $0.45 \times 3.664 =
\qty{1.65}{\celsius}$ प्रति \qty{1000}{PgC}, और परास 0.989 से 2.31 हो जाता है। (ख) $2390 \times
12.011/44.009 = \qty{652}{PgC}$। (ग) $0.45 \times 2.39 = \qty{1.08}{\celsius}$।
\end{solution}

\begin{solution}{exo:b2:anthropogenic-climate-change:5}
(क) $g = 159/398 = 0.400$; $T = (398/\num{5.670e-8})^{1/4} = \qty{289}{K}$। (ख) $-0.6 \times 4
\times \num{5.670e-8} \times 288^3 = \qty{-3.25}{W.m^{-2}.K^{-1}}$; \qty{298}{K} पर, $-3.60$।
(ग) कम संवेदनशील: अनुक्रिया $T^3$ के रूप में बढ़ती है, अतः गर्म ग्रह हर अतिरिक्त अंश को अधिक
आसानी से त्यागता है, अन्य प्रतिपुष्टियाँ समान रहने पर।
\end{solution}

\begin{solution}{exo:b2:anthropogenic-climate-change:6}
(क) $1.85 - 0.50 = \qty{1.35}{W.m^{-2}.\celsius^{-1}}$। (ख) $\sqrt{0.32^2 + 0.20^2} = 0.38$।
(ग) योग किसी भी पद से कहीं बेहतर ज्ञात है: उनकी त्रुटियाँ प्रबल रूप से प्रतिसहसंबद्ध हैं। जहाँ
ऊपरी क्षोभमंडल अधिक गर्म होता है, वहाँ वह अधिक वाष्प भी रखता है, अतः अधिक प्रबल ऋणात्मक \omterm{def:b2:anthropogenic-climate-change:feedbacks}{तापह्रास दर
प्रतिपुष्टि} अधिक प्रबल धनात्मक \omterm{def:b2:anthropogenic-climate-change:feedbacks}{जल-वाष्प प्रतिपुष्टि} के साथ आती है।
\end{solution}

\begin{solution}{exo:b2:anthropogenic-climate-change:7}
(क) $\lambda = (0.59 - 2.20)/1.03 = \qty{-1.56}{W.m^{-2}.\celsius^{-1}}$; $3.93/1.56 =
\qty{2.52}{\celsius}$। (ख) $3.93/(1.56 - 0.5) = 3.93/1.06 = \qty{3.71}{\celsius}$। (ग) $3.93/1.46 =
\qty{2.69}{\celsius}$। मूल्यांकन, पूर्ण वितरणों के साथ, 3.5 और \qty{2.7}{\celsius} पाता है।
\end{solution}

\begin{solution}{exo:b2:anthropogenic-climate-change:8}
(क) $0.617/3682 = \qty{1.68e-4}{\celsius^{-1}}$। (ख) $0.113 \times 0.396 = \qty{0.0447}{m}$, यानी
\qty{44.7}{mm}, बजट के \qty{47.5}{mm} के निकट। (ग) $\num{2e-4} \times 0.2 \times 700 =
\qty{0.028}{m}$, \qty{28}{mm}।
\end{solution}

\begin{solution}{exo:b2:anthropogenic-climate-change:9}
(क) $\sum xy = 0.311$, $\sum x^2 = 0.300$: $\beta = 1.04$। (ख) \qty{95}{\percent} अंतराल
$1.04 \pm 1.96 \times 0.08$ है, 0.88 से 1.20: वह 0 को बाहर रखता है, अतः अनुक्रिया संसूचित है, और
1 को समाहित करता है, अतः प्रेक्षित परिवर्तन अनुकरित परिवर्तन से संगत है। (ग) अंतराल 0.34 से 0.66 अब भी
0 को बाहर रखता है (संसूचित), पर 1 को भी बाहर रखता है: मॉडल अनुक्रिया को लगभग दो के गुणक से
अधिक आँकता।
\end{solution}

\begin{solution}{exo:b2:anthropogenic-climate-change:10}
(क) $\sum g^2 = 1.20$, $\sum a^2 = 0.18$, $\sum ga = -0.46$, $\sum gy = 0.74$, $\sum ay = -0.28$।
प्रसामान्य समीकरण $1.20\,\beta_g - 0.46\,\beta_a = 0.74$ और $-0.46\,\beta_g + 0.18\,\beta_a
= -0.28$ हैं; उनका सारणिक $1.20 \times 0.18 - 0.46^2 = 0.0044$ है, और $\beta_g = (0.74 \times
0.18 - 0.46 \times 0.28)/0.0044 = 1.00$, $\beta_a = (-1.20 \times 0.28 + 0.46 \times 0.74)/0.0044 =
1.00$। (ख) ग्रीनहाउस गैसें $+\qty{0.80}{\celsius}$, एरोसॉल $-\qty{0.30}{\celsius}$। (ग)
सारणिक छोटा है क्योंकि $g$ और $a$ लगभग समानुपाती हैं (उनका सहसंबंध $-0.99$ है): $y$ में थोड़ा
शोर दोनों मापन-गुणकों को साथ-साथ बढ़ाकर समाहित किया जा सकता है, अतः हर एक कम
बाधित रहता है, यद्यपि उनका संयुक्त प्रभाव भली-भाँति ज्ञात है।
\end{solution}

\begin{solution}{exo:b2:anthropogenic-climate-change:11}
(क) $(0.93 - 0.1)/0.45 \times 1000 = \qty{1844}{Gt}$; उसमें से 210 और $26 \times 0.93 = 24$ घटाने पर:
\qty{1610}{Gt}। (ख) $0.93 - x = 0.45 \times (1350 + 210 + 24)/1000 = 0.713$, अतः $x =
\qty{0.22}{\celsius}$: आँके गए 0.1 से \qty{0.2}{\celsius} के लगभग ऊपरी छोर पर। (ग)
बजट गैर-CO$_2$ तापन $x$ पर केवल $-x/\text{TCRE}$ के माध्यम से निर्भर है: हर
\qty{0.1}{\celsius} की कीमत $0.1/0.45 \times 1000 = \qty{222}{Gt}$ है, सीमा चाहे जो हो।
\end{solution}

\begin{solution}{exo:b2:anthropogenic-climate-change:12}
(क) $3.93/1.81 = \qty{2.17}{\celsius}$, $3.93/0.51 = \qty{7.71}{\celsius}$, $3.93/1.16 =
\qty{3.39}{\celsius}$। (ख) निचली ओर $3.39 - 2.17 = \qty{1.22}{\celsius}$ लंबी है, ऊपरी ओर
$7.71 - 3.39 = \qty{4.32}{\celsius}$: फलन $1/|\lambda|$ उत्तल है, अतः शून्य की ओर
$\lambda$ के बराबर क़दम ECS को अधिकाधिक बढ़ाते हैं। (ग) दुर्बल रूप से ऋणात्मक $\lambda$ बहुत
बड़ी ECS देता है, और शून्य के निकट $\lambda$ के छोटे परिवर्तनों को ख़ारिज करना कठिन है; कम ECS के लिए प्रबल
ऋणात्मक $\lambda$ चाहिए, जिसकी सीमा \omterm{def:b2:anthropogenic-climate-change:feedbacks}{प्लांक अनुक्रिया} बाँधती है और प्रेक्षित तापन जिसका खंडन करता है।
\end{solution}

\begin{solution}{pb:b2:anthropogenic-climate-change:1}
\textbf{1.} $2.16 + 0.54 + 0.21 + 0.41 = \qty{3.32}{W.m^{-2}}$।
\textbf{2.} $-0.22 - 0.84 = \qty{-1.06}{W.m^{-2}}$।
\textbf{3.} $3.32 + 0.47 + 0.05 - 1.06 - 0.20 + 0.08 + 0.06 = \qty{2.72}{W.m^{-2}}$, आँका गया
कुल।
\textbf{4.} $3.93\log_2(410/278) = \qty{2.20}{W.m^{-2}}$, 2.16 के \qty{2}{\percent} के भीतर।
\textbf{5.} $3.93\log_2(427.35/278) = \qty{2.44}{W.m^{-2}}$।
\textbf{6.} एरोसॉल--मेघ पद, $-1.45$ से $-0.25$, \qty{1.20}{W.m^{-2}} चौड़ा, कुल के 1.52 के
सामने: वह कुल की अधिकांश अनिश्चितता वहन करता है क्योंकि वह बड़ा है और कम
बाधित है।
\textbf{7.} $0.79 \times \num{5.10e14} \times 12 \times \num{3.156e7} = \qty{1.53e23}{J}$,
\qty{153}{ZJ}।
\textbf{8.} $0.91 \times 153 = \qty{139}{ZJ}$।
\textbf{9.} $\lambda = (0.59 - 2.20)/1.03 = \qty{-1.56}{W.m^{-2}.\celsius^{-1}}$।
\textbf{10.} $3.93/1.56 = \qty{2.52}{\celsius}$।
\textbf{11.} $3.93 \times 1.03/2.20 = \qty{1.84}{\celsius}$।
\textbf{12.} $3.93/1.16 = \qty{3.39}{\celsius}$, \qty{0.87}{\celsius} अधिक।
\textbf{13.} अब तक का तापन वहाँ सबसे प्रबल रहा है जहाँ प्रतिपुष्टियाँ सबसे अधिक ऋणात्मक हैं; जैसे-जैसे
प्रतिरूप विकसित होता है, प्रतिपुष्टि प्राचल कम ऋणात्मक होता जाता है, अतः दीर्घकालिक संवेदनशीलता
प्रभावी संवेदनशीलता से अधिक होती है।
\textbf{14.} $0.45 \times 44.009/12.011 = \qty{1.65}{\celsius}$ प्रति \qty{1000}{PgC}।
\textbf{15.} 0.989 से \qty{2.31}{\celsius} प्रति \qty{1000}{PgC}।
\textbf{16.} $2390 \times 12.011/44.009 = \qty{652}{PgC}$।
\textbf{17.} $0.45 \times 2.39 = \qty{1.08}{\celsius}$।
\textbf{18.} पूर्णांकन के भीतर, अकेली कार्बन डाइऑक्साइड पूरे मानवीय तापन की व्याख्या करती है: अन्य
ग्रीनहाउस गैसों का तापन और एरोसॉलों का शीतलन लगभग एक-दूसरे को निरस्त करते हैं।
\textbf{19.} $1.5 - 1.07 = \qty{0.43}{\celsius}$।
\textbf{20.} $(0.43 - 0.1)/0.45 \times 1000 = \qty{733}{Gt}$ कार्बन डाइऑक्साइड।
\textbf{21.} $733 - 210 = \qty{523}{Gt}$।
\textbf{22.} $26 \times 0.43 = \qty{11.2}{Gt}$; $523 - 11.2 = \qty{512}{Gt}$।
\textbf{23.} वह $0.1/0.45 \times 1000 = \qty{222}{Gt}$ घटकर \qty{290}{Gt} रह जाता है: वही
\qty{220}{Gt} जितना, मूल्यांकन के अनुसार, अन्य चालकों के बारे में चुनाव बजट को खिसकाते हैं।
\textbf{24.} TCRE से परिकलित, 2020 के आरंभ से \qty{1.5}{\celsius} के लिए \omterm{def:b2:anthropogenic-climate-change:budget}{शेष कार्बन बजट}
\textbf{\qty[text-series-to-math=true]{512}{Gt}} कार्बन डाइऑक्साइड है, जो \qty{3}{\percent} के भीतर है उन
\qty{500}{Gt} के, जो \qty{50}{\percent} प्रायिकता पर आँके गए हैं।
\end{solution}
